%%%%%%%%%%%%%%%%%%%%%%%%%%%%%%%%%%%%%%%%%%%%%%%%%%%%%%%%%%%%%%%%%%%%%%%%%%%%%
%集録提出用のサンプルファイルです。
%このファイルを参考にして集録を作成してください。
%%%%%%%%%%%%%%%%%%%%%%%%%%%%%%%%%%%%%%%%%%%%%%%%%%%%%%%%%%%%%%%%%%%%%%%%%%%%%
%%%%%%%%%%%%%%%%%%%%%%%%%%%%%%%%%%%%%%%%%%%%%%%%%%%%%%%%%%%%%%%%%%%%%%%%%%%%%
\documentclass[a4paper,10pt,oneside,twocolumn,notitlepage,final]{jarticle}
\usepackage{ss17(UTF-8)}
%%%%%%%%%%%%%%%%%%%%%%%%%%%%%%%%%%%%%%%%%%%%%%%%%%%%%%%%%%%%%%%%%%%%%%%%%%%%%
%%ss17.styファイルはtexファイルと同じディレクトリに置いてください。
%%ヘッダ、フッターなどの全体のレイアウトは変更しないでください。
%%ここより上は変更しないでください。
%%%%%%%%%%%%%%%%%%%%%%%%%%%%%%%%%%%%%%%%%%%%%%%%%%%%%%%%%%%%%%%%%%%%%%%%%%%%%
%%使いたいパッケージがある場合は以下に書いてください。
\usepackage[dvipdfmx]{graphicx}% 


%%%%%%%%%%%%%%%%%%%%%%%%%%%%%%%%%%%%%%%%%%%%%%%%%%%%%%%%%%%%%%%%%%%%%%%%%%%%%
%%名前、所属、タイトルは以下に記入してください。
%%所属は以下のように（）内にお願いします。
\author{磯谷 和秀 (名古屋大学大学院 理学研究科)}
\title{巨大衝突ステージにおける衝突破壊の重要性:\\
$N$体計算$\cdot$統計的手法のハイブリッドコードの開発}
%%%%%%%%%%%%%%%%%%%%%%%%%%%%%%%%%%%%%%%%%%%%%%%%%%%%%%%%%%%%%%%%%%%%%%%%%%%%%


\begin{document}


%%概要は\abst内に記入してください。
%%\maketitleは必要ありません。
%%以下の\abst{}に概要を記入することにより
%%タイトル、名前、日付、概要が一括して出力されます。
\abst{
太陽系の地球型惑星は、最終段階で火星サイズの原始惑星同士が衝突合体を繰り返し形成される。この巨大衝突ステージにおいて地球や地球-月系が形成される。
一方、太陽系外で起こる巨大衝突ステージは、衝突に伴い放出される破片によりデブリ円盤が形成され、観測されている暖かいデブリ円盤(すなわち地球形成領域のデブリ円盤)を説明することができる[1]。
巨大衝突ステージに形成されるデブリ円盤について調べるためには、原始惑星の長期的軌道進化と、破壊を扱うことができる計算が必要である。
しかし衝突により放出される破片の数は$10^{35}$個以上にもなり、$N$体計算ではとても扱うことはできない。
このような多数の粒子を取り扱うには、一つ一つの粒子を取り扱うのではなく、統計力学に基づいた統計的手法が有効であるが、統計的手法では、破片が重力的に集積する際にサイズ分布が非軸対称になることや、原始惑星による軌道共鳴のような、重力相互作用の取り扱いができない。
すなわち$N$体計算と統計的手法を同時に用いると、軌道進化と破壊を同時に考慮した計算を行うことができる。
そこで本研究では、$N$体計算と統計的手法を組み合わせた、衝突破壊を扱うことができるハイブリッドコードの開発を行う。
多数の破片を少数のトレーサーと呼ばれるスーパー粒子に近似することで$N$体計算のコストを抑える。
またそれぞれのトレーサーの周りに扇形領域[2]を考え、その領域に入った他のトレーサーを用いて表面数密度と平均相対速度を計算し、破壊による天体の減少[3]を取り扱う。
さらに本講演では、ハイブリッドコードにより得られる、巨大衝突ステージにおけるデブリ円盤の明るさの空間分布進化についても議論する。
}



\section{Introduction}
セクションの名前は章立ての目安ですので適宜変更してください。


\section{Methods/Instruments \\ and Observations}
本文を記入してください。本文を記入してください。本文を記入してください。本文を記入してください。本文を記入してください。本文を記入してください。


\section{Results}
本文を記入してください。本文を記入してください。本文を記入してください。本文を記入してください。本文を記入してください。本文を記入してください。


\section{Discussion}
本文を記入してください。本文を記入してください。本文を記入してください。本文を記入してください。本文を記入してください。本文を記入してください。


\section{Conclusion}
本文を記入してください。本文を記入してください。本文を記入してください。本文を記入してください。本文を記入してください。本文を記入してください。


\section{図の挿入方法}
図の挿入方法の簡単な説明です。
figure環境で図を挿入します。
[htbp]で図を入れる位置を指定します。
例のように同時に書いている場合は
h:その場所、t:ページ上部、b:ページ下部、p:独立したページ
の順に挿入しようとします。


%\begin{figure}[htbp]
% \centering
% \includegraphics[width=3cm,clip]{pic.eps}
% \caption{図の挿入方法の説明}
% \label{pic} %ラベルをつけることにより参照できるようになります。
%\end{figure}

\verb+\+includegraphics[width=3cm,clip]\{pic.eps\}
で同じディレクトリにある"pic.eps"という画像ファイルを指定します。
オプションで画像の大きさを指定します。

\verb+\+figref\{ラベル\}でこのように参照できます（例: \figref{pic}）。

\section{表の挿入方法}
表の挿入方法の簡単な説明です。
table環境とtabular環境で表を挿入します。
table環境のオプションで表の位置を決めます。
この方法は図の位置と同じです。
table環境の後に列を決めます。
\{lcr\}のようにすれば一列目は左寄せ、二列目はセンタリング、三列目は右寄せになります。
サンプルでは\{l|c|r\}のようにして列に罫線を入れています。
\verb+\+hline を入れることにより行の罫線を入れることができます。

\begin{table}[htb]
 \caption{表の挿入方法の説明}
  \centering
   \begin{tabular}{|l|c|r|} \hline
     A & B & C  \\ \hline
     1A & 1B & 1C  \\
     2A & 2B & 2C  \\ \hline
   \end{tabular}
  \label{table}%ラベルをつけることにより参照できるようになります。
\end{table}

\verb+\+tabref\{ラベル\}でこのように参照できます（例: \tabref{pic}）。

\section{参考文献の書き方}
\noindent
著者が一人の場合\citep{ラベル1}。\\
著者が二人の場合\citep{ラベル2}。\\
著者が三人の場合\citep{ラベル3}。\\
\section*{Acknowledgement}
謝辞がある場合は記入してください。

\small
\begin{thebibliography}{99}
\bibitem[著者A(2013)]{ラベル1}
 著者A 2013, 発行元1
 \bibitem[著者B \& 著者C(2014)]{ラベル2}
 著者B, \& 著者C 2014,  発行元2
\bibitem[著者D et al.(2015)]{ラベル3}
 著者D, 著者E, \& 著者F 2015, 発行元3
 
 
\end{thebibliography}






\end{document}
